\subsection*{先数清究竟重复了什么}

Alice 的两次点击产生了两次 HTTP 受理、两笔数据库事务和两条扣款流水。讨论幂等性，需要分别指出准备约束其中哪一个对象。

一次请求尝试，是为获取或推进结果而发起的一次通信。R1 与 R2 有各自的连接、开始时间、错误或响应，服务也可以给每次受理分配不同的 \code{request\_id}。

业务操作对应 Alice 授权的一次购买。她希望这次购买成功或得到明确拒绝，为完成它可能发送多次请求。她决定再买一包相纸，则产生另一次业务操作，即便商品、数量与金额完全相同。

业务效果，是操作在选定业务范围中已经成立的变化。本章成功购买的效果包括余额减少 3000、建立一张已付款订单和一条扣款流水，它们在同一事务中提交。事务尝试回滚以后又成功，可能执行了多遍 SQL，最终只有一笔已提交购买。

返回结果说明这笔购买得到了什么处理：成功时包括 P1、O1、D1 和金额 3000，余额不足时说明 P1 以该结果结束。Alice 还可以从账户历史中看到订单与流水，其他已授权设备也可以查询同一笔购买。因此，检查购买是否正确完成，需要同时核对扣款、记录和这些查询取得的结果。

\begin{figure}[htbp]
\centering
\begin{tikzpicture}[x=1cm,y=1cm]
\node[booknode,text width=5cm] (u) at (5,0) {\mbox{业务操作 $u$}\\\mbox{Alice 决定买一包相纸}};
\node[booknode,text width=2.5cm] (r1) at (1.5,-1.8) {请求尝试 R1};
\node[booknode,text width=2.5cm] (r2) at (8.5,-1.8) {请求尝试 R2};
\draw[bookarrow,dashed] (u)--node[left,font=\scriptsize]{首次发送}(r1);
\draw[bookarrow,dashed] (u)--node[right,font=\scriptsize]{为完成原购买重发}(r2);
\node[booknode,text width=3.7cm] (p1) at (1.5,-3.5) {当前入口创建 P1\\提交一笔 3000 扣款};
\node[bookresult,text width=3.7cm] (p2) at (8.5,-3.5) {当前入口创建 P2\\再提交一笔 3000 扣款};
\draw[bookarrow] (r1)--(p1);
\draw[bookalert] (r2)--(p2);
\node[bookresult,text width=7.8cm] at (5,-5.1)
 {按同一业务操作统计：$E(u)=2$\\所需上限：$E(u)\leq 1$};
\node[booknote,text width=11cm] at (5,-6.1)
 {虚线按 Alice 的购买目的归类；当前数据库尚未保存这项跨请求关系};
\end{tikzpicture}
\caption{一次操作可以经历多次通信尝试。效果数量应围绕业务操作统计。}
\label{fig:ch02-operation-model}
\end{figure}


把同一业务操作记为 $u$，把某个观察点以前属于 $u$ 的已提交成功扣款条数记为 $E(u)$。最多一次效果可以写成：
\[
 E(u)\leq 1.
\]

如果这是一笔有效、可以支付且没有被合法取消的购买，系统恢复可用后又持续执行恢复流程，还希望最终达到 $E(u)=1$，就需要另行写出这些进展条件。余额不足是允许的终态，此时 $E(u)=0$。

当前数据库还没有这个跨请求身份 $u$。对 P1 数一次，再对 P2 数一次，得到的是每笔数据库购买各自的记录数。判断 Alice 的这次意图只扣一次，首先需要把 R1、R2 归入同一操作，再检查它们的所有生效路径。

\subsection*{先选定业务观察的范围}

只研究余额时，一个状态可以是数值 10000；研究整笔购买时，还要包含购买、订单和流水。从 10000 扣 3000，一次得到 7000，两次得到 4000，单看余额已经能发现重复支出。若换成每次都把余额设置为固定值 7000，紧邻执行两次，余额会保持 7000。要判断购买是否只产生一份效果，还得继续看其他记录。

继续看设置余额为 7000。若每次设置都新增一条 3000 分扣款流水，第一次以后是余额 7000、一条流水；第二次以后是余额 7000、两条流水。余额相同，完整业务状态却已变化。

\begin{figure}[htbp]
\centering
\begin{tikzpicture}[x=1cm,y=1cm]
\node[font=\small\sffamily\bfseries] at (2.5,.7) {第一次设置余额并记流水};
\node[font=\small\sffamily\bfseries] at (10,.7) {第二次重复处理};
\node[booknode,text width=4.7cm] (a) at (2.5,-.2) {余额：7000};
\node[booknode,text width=4.7cm] (b) at (10,-.2) {余额：7000};
\draw[bookarrow] (a)--node[above,font=\scriptsize]{只看余额：相同}(b);
\node[booknode,text width=4.7cm] (c) at (2.5,-1.8) {扣款流水\\D1：3000};
\node[bookresult,text width=4.7cm] (d) at (10,-1.8) {扣款流水\\D1：3000\\D2：3000};
\draw[bookalert] (c)--node[above,font=\scriptsize]{流水多一条}(d);
\draw[Ink!25,dashed] (-.2,-3.1)--(12.7,-3.1);
\node[booknote,text width=5cm] at (2.5,-3.8) {完整状态\\余额与一条扣款记录};
\node[booknote,text width=5cm] at (10,-3.8) {完整状态\\余额与两条扣款记录};
\end{tikzpicture}
\caption{只比较余额会漏掉新增流水。状态集合的选择决定幂等等式实际检查了哪些变化。}
\label{fig:ch02-state-projection}
\end{figure}


购买与订单也一样：每次新增随机编号的记录，就会改变这些记录集合。讨论幂等性，首先要说明新建对象怎样保持同一身份，以及哪些变化纳入效果范围。固定赋值在其他交易插入后的后果，以及数学幂等等式的准确含义，将在第 4.5 节结合完整状态推演。

两条流水各写 3000，余额却只少 3000，会破坏第一章的核对关系；余额少 6000，流水却只有 3000，也不符合核对规则。完整购买需要把必要记录与余额变化一起约束。

HTTP 幂等方法的定义选择了另一种明确范围：用户请求的预期效果。RFC 9110 允许服务为每次请求分别记日志，也允许重试响应不同\cite{http}。重复 GET 留下多行访问日志，与给 Alice 额外扣一次款，分别属于传输审计和购买效果。传输审计记录每次通信尝试发生了什么，购买效果则描述账户与订单实际发生的变化。

设计时可以保留每次尝试的审计日志，扣款流水则逐笔对应真实扣款。购买成功通知是否允许重复，需要通知入口自己的处理规则。以后加入邮件、外部付款或消息系统时，还需要逐一检查它们的效果。

\subsection*{效果只成立一次，返回结果仍可能出错}

考虑另一条独立历史：Alice 初始余额为 2000 分，P1 购买 3000 分相纸，保存余额不足终态，响应随后丢失。后来充值使余额足够，恢复 P1 仍应读回原拒绝；充值后的重新购买需要另一项确认。第 5.4 节将完整给出这两类响应的字段与处理规则。

即使某种机制挡住了第二次扣款，原结果仍需要保存。余额 7000 可能来自购买成功，也可能来自人工调整；当前余额不足也可能由后来的其他消费造成。从现值还原原购买，会把这些历史混在一起。

前面提到，Raft 让多个副本确定一致的已提交命令顺序。在相同初始数据和确定的处理规则下，各副本依次执行这些命令，就能得到相同状态。这种由当前状态和输入命令决定下一状态及结果的模型，称为状态机；购买服务可以把扣款看作命令，把余额和购买记录看作状态。

Diego Ongaro 的博士论文用锁服务说明返回结果的问题。锁服务记录某项资源当前由谁占用；客户端获取锁成功以后，可以执行需要独占这项资源的动作，其他调用者要等锁释放后再获取。客户端第一次获取已经成功，确认却丢失了；重试再次执行获取动作，看到锁已被占用，于是得到获取失败\cite{ongaro}。锁仍只有一把，拥有者也没变，返回结果却让原调用者误以为自己没有取得锁。

\begin{figure}[htbp]
\centering
\begin{tikzpicture}[x=1cm,y=1cm]
\foreach \x in {0,6,11.5} {\draw[Ink!35,dashed] (\x,-.5)--(\x,-5.8);}
\node[booknode,text width=2cm] at (0,0) {客户端 A};
\node[booknode,text width=2cm] at (6,0) {锁服务};
\node[booknode,text width=2cm] at (11.5,0) {锁状态};
\draw[bookarrow] (0,-1)--node[above,font=\footnotesize]{获取锁}(6,-1);
\draw[bookarrow] (6,-1.7)--(11.5,-1.7);
\node[booknode,text width=2.3cm] at (11.5,-2.4) {已由 A 持有};
\draw[bookalert,dashed] (6,-2.8)--node[above,font=\footnotesize]{成功确认丢失}(1,-2.8);
\draw[bookarrow] (0,-3.5)--node[above,font=\footnotesize]{重试获取锁}(6,-3.5);
\draw[bookalert] (6,-4.3)--node[above,font=\footnotesize]{重新执行：锁被占有，返回失败}(0,-4.3);
\node[bookresult,text width=3.5cm] at (0,-5.4) {误以为未取得锁};
\node[booknode,text width=2.3cm] at (11.5,-5.4) {仍由 A 持有};
\end{tikzpicture}
\caption{最终锁状态保持不变，重试结果却误导调用方。保存并返回原结果也是重复调用处理的一部分。}
\label{fig:ch02-lock-result}
\end{figure}


对购买也需要同时检查效果与原结果。原扣款成立后，服务应有与这笔操作关联的完成信息，重试才能取回正确结论。

RIFL 把这份信息叫作完成记录，要求它包含 RPC 标识和返回给客户端的结果，并与业务变化原子创建、可靠保存\cite{rifl}。原子创建指两项修改共同成立：业务生效时也有完成记录，整组未完成时两项都不生效。这样把操作事实和结果一起保留下来，恢复者才有依据回答原调用。后面的购买协议会具体安排这项写入，以及重复请求怎样读回原记录并重新返回结果；这种重新提供已保存结果的处理称为结果重放。

\subsection*{结果稳定不要求每个响应字节相同}

P1 已经成功，后来查询它，仍应说明商品 \code{paper-001}、金额 3000、订单 O1 和扣款 D1。账户随后又消费，改变的是当前余额，P1 的成功结果保持原样。

同一查询在不同时间返回的 \code{Date}、\code{request\_id} 和链路追踪信息可以不同。\code{Date} 是 HTTP 响应头中说明消息生成时间的字段；链路追踪信息用于把一次请求经过各服务的处理步骤关联起来，便于排查本次通信。它们描述的是当前响应。首次 POST 返回 201，详情 GET 返回 200，也分别符合这两个入口的用途。

动态字段需要说明时间含义。\code{balance\_after\_minor=7000} 描述 D1 当次扣款后的余额，之后仍是 7000；另一笔独立购买以后，当前账户余额可以是 4000。页面可以同时显示这两个值，各自对应不同的记录与时刻。

\begin{figure}[htbp]
\centering
\begin{tikzpicture}[x=1cm,y=1cm]
\node[font=\small\sffamily\bfseries] at (2.5,1.1) {P1 的稳定业务结果};
\node[font=\small\sffamily\bfseries] at (10,1.1) {查询时的当前信息};
\node[booknode,text width=5.1cm] (fixed) at (2.5,-.7)
 {P1 成功，商品 paper-001\\金额 3000，订单 O1，扣款 D1\\D1 当次扣后余额 7000};
\node[booknode,text width=4.3cm] (now1) at (10,0) {第一次查询\\当前余额 7000};
\node[booknode,text width=4.3cm] (now2) at (10,-1.8) {另一笔购买以后再查询\\当前余额 4000};
\draw[bookarrow] (now1)--node[right,font=\scriptsize]{账户继续变化}(now2);
\node[booknode,text width=11cm] at (6.1,-3.4)
 {每次响应另有自己的 Date、request\_id 与追踪信息\\业务结果可保持稳定，查询元数据随本次通信产生};
\draw[bookarrow,dashed] (fixed.east)--node[above,font=\scriptsize]{一并返回}(now1.west);
\end{tikzpicture}
\caption{当次扣后余额属于原流水，当前余额属于查询时刻。区分字段的时间含义，才能解释同一购买的不同响应。}
\label{fig:ch02-result-time}
\end{figure}


因此可以在返回内容中区分原操作的业务结果与本次响应的信息。若接口进一步承诺精确回放原响应体，就需要保存相应内容和表示版本；表示版本说明当时使用哪一版字段及编码规则，便于接口升级后仍按原样返回。只要求业务结果稳定时，则可让本次查询的时间与追踪信息正常更新。

\subsection*{从其他观察者检查同一条历史}

Alice 的手机没收到 P1，但另一台已授权设备可能已经在账户历史中读到 P1，看到余额 7000。假设服务处理重试时删除 P1，再创建 P2，手机最后取得了成功；另一台设备随后查询 P1，却得到不存在。

P1 曾经作为成功购买被查询过，删除以后，两台设备取得的信息就失去对应。恢复处理需要沿用已保存的购买事实，使此前取得 P1 的调用方仍能继续查询这笔购买。

\begin{figure}[htbp]
\centering
\begin{tikzpicture}[x=1cm,y=1cm]
\node[booknode,text width=2.7cm] (phone) at (0,0) {Alice 的手机\\P1 响应丢失};
\node[booknode,text width=3.3cm] (p1) at (5,0) {已提交的 P1\\O1、D1，余额 7000};
\node[booknode,text width=2.7cm] (other) at (10.3,0) {另一台授权设备\\已经读到 P1};
\draw[bookarrow] (p1)--(other);
\node[bookresult,text width=3.3cm] (erase) at (5,-1.9) {若删除 P1 并改建 P2};
\draw[bookalert] (p1)--(erase);
\node[bookresult,text width=2.7cm] (missing) at (10.3,-1.9) {再次查询 P1\\记录消失};
\draw[bookalert] (erase)--(missing);
\node[booknote,text width=2.7cm] at (0,-1.9) {手机最终得到成功\\还需核对\\完整业务历史};
\draw[Ink!25,dashed] (-1.5,-3)--(11.8,-3);
\node[booknode,text width=4.7cm] (debit) at (1,-4) {保留原购买与扣款 D1};
\node[booknode,text width=4.7cm] (refund) at (8.9,-4) {若需返还，新增退款 F1\\引用 D1，记录返还金额};
\draw[bookarrow] (debit)--(refund);
\end{tikzpicture}
\caption{其他观察者已经取得的记录也需要延续。补偿通过新增关联记录修正经济结果，保留此前的扣款历史。}
\label{fig:ch02-observers}
\end{figure}


发现两次扣款后，直接删除 D2 并把余额加回 3000，也会让已读取 D2 的对账程序丢失历史依据。退款可以新增返还记录，引用原扣款。它修正经济结果，同时保留原购买、扣款和返还的先后关系。

本章的账户、购买、订单和扣款流水，都在同一数据库的一笔事务里修改；这个能够共同提交或回滚的范围，称为事务域。如果成功订单同时触发了外部付款或通知，本地回滚的范围就止于本地数据库，外部动作需要各自的恢复办法。后面扩展执行者时，仍会检查每个参与者取得哪些事实，以及恢复怎样延续这些已经发生的变化。

\subsection*{重新核对开头的两次扣款}

R1 没有返回业务结果，R2 返回成功；数据库中两笔事务都完整提交；账户累计减少 6000，存在两条流水；Alice 只授权了买一包相纸，其中一笔因此属于重复生效。这些描述分别统计了请求、事务和业务效果。

两笔事务满足原子性，余额 4000 满足非负约束，成功响应也只出现了一次。要指出这次重复扣款，需要检查同一购买意图最多产生一条成功扣款。还要同时允许两次真实独立购买分别成立。

一个改进版本应当回答几个具体问题：同一操作重试以后，余额是否再次减少；原结果在进程重启后怎样取回；从未生效的请求在系统恢复后怎样继续完成；两次独立购买怎样分别保留。前面每条失败历史，都给出了其中一项要求的来由。

RIFL 的调用标识与 Raft 的客户端序号，把多次尝试与同一命令联系起来。客户端序号为每个客户端发起的新命令依次编号，重试沿用原编号；服务器再由保存的编号与结果处理重复调用\cite{rifl,raft}。购买接口也需要先建立这样的关系。

Alice 第一次点击时，还没有从响应取得 P1。因此，操作身份的生成与保存需要发生在结果可能丢失以前。下一章将具体说明这项身份由谁生成、怎样随请求传递、与账户及输入怎样关联，以及另一次相同商品购买怎样获得新的身份。
\FloatBarrier
